\section{问题四：波动电价下的情景采购与连续储能调度}
\label{sec:q4}
\subsection{信息条件与购电权限}
波动电价下，仅按点预测购电可能低估供电缺口，而紧急电价为正常电价的5倍，缺口造成的费用与余电损失并不对称。因此，先建立点预测滚动方案作为基准，再利用已实现的负载、光伏与价格联合误差构造采购情景，在正常购电和缺口风险之间作出权衡。

本节采用2025年1月作为点预测的历史初始化期，自2月1日以6\,000 kWh连续运行至2026年1月1日，储电量不按日重置。区间实际负载、光伏和电价在区间结束后进入历史；计划中的未来量由当时可用数据预测。情景误差库自2月1日空初始化，前七个完整运行日采用原点预测方案，随后使用已经实现的预测误差。

将题目要求的两类条件分别记为问题四（2）和问题四（3）。本文重点报告问题四（2）的每日冻结合同结果：每日0:00提交144个区间的初始合同$G_k^0$，当日不再调整；问题四（3）增加已发布光伏预报，并在6、12、18时调整当日尚未开始区间的合同，作为同一情景机制下的对照。跨日窗口内的次日购电量只用于预期调度，次日0:00才形成实际合同。

整体求解流程见图\ref{fig:q4-workflow}。发布时更新点预测与情景，执行时只采用当前共同动作；实际状态和已实现误差分别传递至后续调度和情景学习。
\begin{figure}[H]
\centering
\PaperGraphic[width=0.98\linewidth]{figures/q4_scenario_workflow.pdf}
\caption{问题四的情景采购与实际反馈流程}
\label{fig:q4-workflow}
\end{figure}

\subsection{点预测与滚动窗口}
\subsubsection{历史基线与残差修正}
设$t$为当前决策时刻，也是最近已结束区间的右端点，$s$为连续十分钟序列的索引。利用历史同刻数据构造负载与价格基线：
\begin{align}
 b_s^{\mathrm L}&=\frac{8P_{s-1008}^{\mathrm L}+4P_{s-2016}^{\mathrm L}+2P_{s-3024}^{\mathrm L}+P_{s-4032}^{\mathrm L}}{15},\\
 b_s^p&=0.5p_{s-144}+0.3p_{s-1008}+0.2p_{s-2016}.
\end{align}
其中144格为一日、1008格为一周。负载基线综合过去四周同刻值，价格基线综合前一日及前两周同刻值。对序列$x$定义残差$e_s^x=x_s-b_s^x$，以截至刷新时刻的有效连续残差对估计
\begin{equation}
 \phi_x=\operatorname{clip}\!\left(\frac{\sum_s e_{s-1}^xe_s^x}{\sum_s(e_{s-1}^x)^2},0,0.999\right),\qquad
 \widehat x_{t+h}=\max\{0,b_{t+h}^x+\phi_x^he_t^x\}.
\end{equation}
这里$h$为未来十分钟步数。衰减项描述近期偏离历史基线的持续性，非负截断用于功率与价格预测。分母为零时取$\phi_x=0$；所需滞后或连续残差不足时停止预测，不跨数据缺口拼接样本。基线权重固定，不依据回放期误差或费用反向调参。

\subsubsection{光伏预报组合与时间细化}
问题四（2）采用过去七日同刻实际光伏均值。问题四（3）对同一目标整点组合全部已发布的官方预报。设版本$v$在发布时的原始提前小时数为$j(v)$，截至刷新时刻$r$的已实现历史误差为$\mathrm{MAE}_{j(v)}(r)$，则版本权重为
\begin{equation}
 w_v(r)=\frac{[\mathrm{MAE}_{j(v)}(r)+1\,\mathrm{kW}]^{-2}}
 {\sum_a[\mathrm{MAE}_{j(a)}(r)+1\,\mathrm{kW}]^{-2}}.
\end{equation}
提前小时数使用原版本属性，历史误差只计已观测目标；误差桶为空时不生成预测。得到每个目标整点的唯一预测后，用线性插值细化为十分钟右端点功率，左端代理取最近结束区间的实际光伏均值，最后乘$1/6$换算为电量。

负载、价格与光伏均在0、6、12、18时刷新，发布间固定预测版本及相关参数。最近刷新时刻为$r$时，时刻$t$的滚动窗口为
\begin{equation}
 \mathcal H_t=[t,\min\{r+24\,\mathrm h,T\}),\qquad T=\text{2026年1月1日0:00}.
\end{equation}
0:00窗口须覆盖完整当日合同；年末窗口截至$T$，仅预测必要目标。点预测和后述情景均在发布间保持不变，窗口随时间推进取剩余部分，不重新划分原始预测时距。

\subsection{联合误差情景的构造}
\subsubsection{已实现误差与时距分组}
设$a$为原预测发布时间，$u$为目标区间右端点，将当时保存的原始点预测与目标实现后的观测值对应，得到联合误差
\begin{equation}
 \boldsymbol\varepsilon_{a,u}=\left(
 P_u^{\mathrm{L,actual}}-\widehat P_{a,u}^{\mathrm L},\,
 P_u^{\mathrm{PV,actual}}-\widehat P_{a,u}^{\mathrm{PV}},\,
 p_u^{\mathrm{actual}}-\widehat p_{a,u}\right).
\end{equation}
仅当$u\le r$且三类实际序列均已观测时，误差行才进入刷新时刻$r$的样本库。按原始时距$u-a$划分$(0,6]$、$(6,12]$、$(12,18]$、$(18,24]$小时四组，每组仅保留$r-28\,\mathrm d<u\le r$的已实现样本。同一目标对应不同发布时间时，各发布--目标对分别计数并等权；不事后回算旧预测。

对第$b$组样本定义净需求误差
\begin{equation}
 \varepsilon_{a,u}^{\mathrm{net}}
 =\frac{\varepsilon_{a,u}^{\mathrm L}-\varepsilon_{a,u}^{\mathrm{PV}}}{6}
 \quad (\mathrm{kWh}).
\end{equation}
计算其经验20\%和80\%分位数$q_{b,0.2},q_{b,0.8}$，分位数采用线性插值。分别选取净需求误差最接近相应分位数的\emph{完整联合误差行}：
\begin{equation}
 \boldsymbol\varepsilon_b^{(\alpha)}
 =\boldsymbol\varepsilon_{a^*,u^*},\qquad
 (a^*,u^*)\in\arg\min_{(a,u)\in\mathcal B_b(r)}
 \left|\varepsilon_{a,u}^{\mathrm{net}}-q_{b,\alpha}\right|,
 \quad\alpha\in\{0.2,0.8\}.
\end{equation}
并列时依次按目标时刻、发布时间取较早样本。三类误差取自同一历史行，以保留该样本中的负载、光伏与价格关系，避免分别取分位数后拼成未出现过的组合。

\subsubsection{三情景输入}
取情景集合$\Omega=\{-,0,+\}$，分别对应低净需求、名义及高净需求情景，令$\boldsymbol\varepsilon_b^-=\boldsymbol\varepsilon_b^{(0.2)}$、$\boldsymbol\varepsilon_b^+=\boldsymbol\varepsilon_b^{(0.8)}$，权重固定为
\begin{equation}
 (\rho_-,\rho_0,\rho_+)=(0.25,0.50,0.25),\qquad
 \boldsymbol\varepsilon_b^0=\boldsymbol0.
\end{equation}
将各时距组的代表误差加至原点预测，并作非负截断：
\begin{align}
 \ell_k^\omega&=\frac{\max\{0,\widehat P_k^{\mathrm L}+\varepsilon_b^{\mathrm{L},\omega}\}}{6},\\
 v_k^\omega&=\frac{\max\{0,\widehat P_k^{\mathrm{PV}}+\varepsilon_b^{\mathrm{PV},\omega}\}}{6},\qquad
 p_k^\omega=\max\{0,\widehat p_k+\varepsilon_b^{p,\omega}\}.
\end{align}
这里$b$由目标相对原发布时间的时距确定，发布间不重选代表样本。冷启动期或空组取零误差；全部情景相同时直接求解原点预测问题。28日窗口、七日冷启动、两处分位数与三情景权重均在候选试验前固定。该构造是对历史联合误差的有限情景近似，不将情景权重解释为经验证的真实发生概率，也不改变原点预测及其误差报表。

\subsection{情景采购与储能优化模型}
\subsubsection{供需、储能及紧急购电约束}
三个情景共享窗口内全部合同$G_k$、交易增减量及充放电量$C_k,D_k$，因此储能状态$E_k$也相同；光伏利用、弃光、合同余电和紧急购电则允许随情景变化。对每个$\omega\in\Omega$建立供需关系：
\begin{align}
 G_k+V_k^\omega+D_k+U_k^\omega&=\ell_k^\omega+C_k+S_k^{\mathrm{grid},\omega},\\
 V_k^\omega+S_k^{\mathrm{PV},\omega}&=v_k^\omega,\qquad
 0\le S_k^{\mathrm{grid},\omega}\le G_k.
\end{align}
各电量非负，共同储能动作满足
\begin{align}
 E_{k+1}&=E_k+0.9C_k-D_k/0.9,\qquad1200\le E_k\le10800,\\
 C_k&\le Mz_k,\qquad D_k\le M(1-z_k),\qquad D_k\le\ell_k^\omega\quad(\forall\omega).
\end{align}
取$M=5000/6$，$z_k\in\{0,1\}$。对每个情景引入紧急购电许可变量$y_k^\omega\in\{0,1\}$，约束
\begin{align}
 U_k^\omega&\le\ell_k^\omega y_k^\omega,\qquad C_k\le M(1-y_k^\omega),\\
 S_k^{\mathrm{grid},\omega}&\le G_k^{\max,\omega}(1-y_k^\omega),\qquad
 S_k^{\mathrm{PV},\omega}\le v_k^\omega(1-y_k^\omega).
\end{align}
紧急电量仅补充负载缺口，不能用于充电，也不能与弃置可用供电同时发生。每个情景沿用原合同上界：新增及次日预期合同取$G_k^{\max,\omega}=\ell_k^\omega+M$，可调整合同取$\max\{G_k^{\mathrm{prev}},\ell_k^\omega+M\}$，冻结合同固定为已确认值。共同$G_k$必须同时满足各情景上界，不为提高可行性放宽购电边界。这些上界为规划模型边界，并非外网物理容量。

全窗口充放电动作共享，使该模型属于受限的情景优化：它允许供需分配随情景变化，但未允许未来电池动作依情景重新决策。因此，不能将所得解称为具有完整未来追索决策的随机规划最优策略。

\subsubsection{合同调整与目标函数}
合同相对上一确认版本的变化写为
\begin{align}
 G_k&=G_k^{\mathrm{prev}}+\delta_k^+-\delta_k^-,\notag\\
 &\delta_k^+\le G_k^{\max,\omega}w_k\quad(\forall\omega),\quad
 \delta_k^-\le G_k^{\mathrm{prev}}(1-w_k),\quad w_k\in\{0,1\}.
\end{align}
由共同方向变量$w_k$限制同时增减购。记$\mathcal N_t$为当前新增及次日预期合同区间集合，$\mathcal A_t$为当前可调整集合，按三情景权重最小化规划费用：
\begin{align}
 \min J_t={}&\sum_{\omega\in\Omega}\rho_\omega\Bigg[
 \sum_{k\in\mathcal N_t}p_k^\omega G_k
 +p_t^\omega\sum_{k\in\mathcal A_t}(1.5\delta_k^++0.5\delta_k^-)\notag\\
 &\hspace{4.5em}+\sum_{k\in\mathcal H_t}5p_k^\omega U_k^\omega\Bigg]
 -v_rE_{\mathrm{end}}.
 \label{eq:q4-scenario-objective}
\end{align}
过去已支付费用为常数，未来调整不提前免费执行。$p_t^\omega$取各情景窗口首区间价格；普通窗口的$v_r$仍取刷新时完整144格\emph{名义}预测价均值的0.9倍，不叠加其他残值候选。模型在各情景均满足设备和合同约束后，权衡新增合同费、当前调整费、情景缺口费与末态储能价值。

\subsubsection{年度末态与储能预留}
窗口触及回放终点$T$时去除末态价值项，改用$E_T=6000$。若距终点剩余$n$步，还须满足必要可达界
\begin{equation}
 6000-0.9nM\le E\le6000+\frac{nM}{0.9},
\end{equation}
并与设备运行范围相交。考虑到冻结合同和预测误差可能限制后续充电，最后24小时在全部状态边界增加预留下界：
\begin{equation}
 E(\tau)\ge6000\ \mathrm{kWh},\qquad T-24\,\mathrm h\le\tau\le T.
\end{equation}
提前跨入该日期的预测窗口也施加此界。该安排减少年末先放电、再依赖有误差预测充回目标的风险，可能增加购电费用；实际末态仍须满足6\,000 kWh等式，不能以预测可达代替实际闭合。

\subsection{实际反馈与费用结算}
每次求解后仅执行当前区间的共同意图$(\overline C,\overline D)$，其余动作留待后续滚动更新。记实际负载、光伏和合同电量为$L,P,G$，根据及时响应假设，本地保护只减少原动作：
\begin{align}
 C&=\min\{\overline C,\max(G+P-L,0)\},\qquad D=\min\{\overline D,L\},\\
 U&=\max\{L+C-G-P-D,0\},\qquad V=\min\{P,L+C-D\},\\
 G_{\mathrm{use}}&=L+C-D-U-V,\qquad
 S^{\mathrm{grid}}=G-G_{\mathrm{use}},\quad S^{\mathrm{PV}}=P-V.
\end{align}
实际供需用于限充、限放及缺口补电，不据此重选经济动作。以实际$C,D$推进储电量，作为下一次滚动优化的初态；非法储能意图或不能及时响应的违规保留，不通过截断储电量修正。

回放期费用按实际合同和交易记录计算：
\begin{align}
 F_{\mathrm{plan}}&=\sum_kp_k^{\mathrm{actual}}G_k^0,\\
 F_{\mathrm{adjust}}&=\sum_vp_v^{\mathrm{trade,actual}}\sum_k(1.5\Delta_{v,k}^++0.5\Delta_{v,k}^-),\\
 F_{\mathrm{emergency}}&=\sum_k5p_k^{\mathrm{actual}}U_k,\\
 F_{\mathrm{total}}&=F_{\mathrm{plan}}+F_{\mathrm{adjust}}+F_{\mathrm{emergency}}.
\end{align}
增减购量取相邻确认版本差额，交易价取该交易时刻所在实际区间价，于区间结束后结算。未用合同仍付费，减购无退款，最终合同不再重复全额计费，滚动窗口目标也不累加为实际总费。问题四（2）的调整费为零。

例如合同依次为$10\to8\to9$ kWh，执行价为2元/kWh，两次调整交易价为1、3元/kWh，紧急购电1 kWh，则总费为$20+5.5+10=35.5$元。该例用于说明逐版结算方式。

\subsection{求解与比较设置}
采用稀疏约束矩阵和SciPy/HiGHS求解，相对间隙阈值为$10^{-4}$，单次限制10秒，超时最多以60秒重试；仅接受求解成功且各情景独立约束复核通过的解。若预测与情景版本、窗口终点、合同和实际储电量与原计划后继条件一致，且不涉及当前新提交或调整合同，可复用可行后缀；扣除已付合同及过去费用后，以原绝对误差界重新核对剩余加权问题的间隙。情景生成或约束不可行时保留失败，不静默改变权重与边界。

各合同条件下，将情景采购与原点预测基准进行单因素比较：实际输入、原负载/光伏/价格预测、合同权限、实际反馈、收费、末态价值及年末边界均相同，仅改变规划中的情景机制。两方案均从同一初态连续运行334日，不分月拼接。另将长窗光伏处理、显式采购余量和普通窗口残值作为独立候选，分别留存完整回放结果，不把未经联合回放的机制叠加成新方案。

记录实际总费及各费用分项、紧急量、弃光、合同余电、保护次数和储能吞吐，并分别评价6小时、24小时点预测误差。以完整可行回放中的实际总费选择本轮方案。该年份已经用于开发诊断，时间因果检验保证决策未读取未来数据，不能代替独立年份验证。

\IfFileExists{tables/q4_results_paper.tex}{% Explicitly bound scenario-procurement results; no model run or approval mutation.
% selection SHA-256: 49d2e7b8422a4cc3942bbea734baa3d219417d89d0dd67d0859cf4f3de17324f
% q4_2: q4-2-v4-scenario-procurement-annual-main-002
% q4_3: q4-3-v4-scenario-procurement-annual-main-002
\subsection{问题四（2）及对照条件的回放结果}
点预测基准与情景采购均从6\,000 kWh初态独立连续运行，覆盖2025年2月至12月的334日、48\,096个区间，实际末态均为6\,000 kWh。表\ref{tab:q4-annual}给出同条件实际费用。
\begin{table}[H]\centering\small
\caption{问题四点预测基准与情景采购的实际费用（单位：万元）}\label{tab:q4-annual}
\begin{tabular}{llrrrr}\toprule
合同条件 & 采购方案 & 初始合同费 & 调整费 & 紧急费 & 总费 \\ \midrule
每日冻结 & 点预测基准 & 1296.3714 & 0.0000 & 421.7606 & 1718.1320 \\
每日冻结 & 情景采购 & 1396.0073 & 0.0000 & 195.9010 & 1591.9083 \\
\midrule
允许调整 & 点预测基准 & 1315.2016 & 45.7007 & 225.7140 & 1586.6164 \\
允许调整 & 情景采购 & 1408.4948 & 30.4363 & 101.3661 & 1540.2971 \\
\bottomrule\end{tabular}\end{table}

\subsubsection{问题四（2）：每日冻结合同}
每日冻结条件下，情景采购比点预测基准节省1\,262\,236.91元（7.3466\%）。初始合同费增加99.6359万元，实际紧急费减少225.8596万元。紧急购电量由83.5854万kWh降至37.8659万kWh，但合同余电增加66.8023万kWh、弃光增加18.7383万kWh。因此，节费来自以较低代价的正常合同替代高价缺口补电，而不是无代价增加供给。

作为对照，允许调整条件下的情景采购比点预测基准节省463\,192.69元（2.9194\%）。初始合同费增加93.2931万元，实际紧急费减少124.3479万元，调整费另减少15.2645万元。

图\ref{fig:q4-procurement}将费用分项和各月费用差并列。初始合同费增加而紧急费减少，净节省由三项账单共同决定。2月每日冻结方案费用比基准高1.8448万元，9月允许调整方案也高1.2803万元，其余月份费用均降低。因此，回放期节费并非逐月一致。
\begin{figure}[H]\centering
\PaperGraphic[width=0.98\linewidth]{figures/q4_procurement_comparison.pdf}
\caption{情景采购与点预测基准的实际费用分解及月费用差}\label{fig:q4-procurement}
\end{figure}

\subsubsection{供电缺口与余电代价}
费用降低伴随更高的合同余电与弃光，表\ref{tab:q4-operation}给出实际运行分项。
\begin{table}[H]\centering\small
\caption{情景采购的实际运行变化（电量单位：万kWh）}\label{tab:q4-operation}
\begin{tabular}{lrrrr}\toprule
 & \multicolumn{2}{c}{每日冻结} & \multicolumn{2}{c}{允许调整} \\
\cmidrule(lr){2-3}\cmidrule(lr){4-5}
指标 & 点预测 & 情景采购 & 点预测 & 情景采购 \\ \midrule
紧急购电量 & 83.5854 & 37.8659 & 50.6744 & 22.1377 \\
合同余电 & 54.3258 & 121.1281 & 76.7139 & 140.5479 \\
弃光电量 & 138.0853 & 156.8237 & 145.6286 & 163.3391 \\
电池吞吐量 & 1266.8205 & 1291.2090 & 1269.2923 & 1285.4029 \\
动作保护次数（次） & 10707 & 5361 & 9016 & 4392 \\
\bottomrule\end{tabular}\end{table}

每日冻结条件下，紧急购电量减少45.7195万kWh，但合同余电增加66.8023万kWh，弃光增加18.7383万kWh。
允许调整条件下，紧急购电量减少28.5367万kWh，但合同余电增加63.8339万kWh，弃光增加17.7105万kWh。
动作保护次数也有所减少，但电池吞吐量增加。当前目标未计寿命损耗，不能据总电费降低推断电池全寿命费用改善。

\subsubsection{点预测误差与候选方案比较}
表\ref{tab:q4-prediction}列出两采购方案共同的点预测误差。情景采购改变了规划中的风险处理，保留原始点预测及其来源，因而费用改善来自采购安排的变化。
\begin{table}[H]\centering\small
\caption{问题四两采购方案共同的点预测误差}\label{tab:q4-prediction}
\begin{tabular}{lrrrr}\toprule
 & \multicolumn{2}{c}{6小时} & \multicolumn{2}{c}{24小时} \\
\cmidrule(lr){2-3}\cmidrule(lr){4-5}
变量 & MAE & RMSE & MAE & RMSE \\ \midrule
负载（kW） & 163.91 & 227.69 & 166.98 & 233.88 \\
光伏（kW，每日冻结） & 153.40 & 303.86 & 153.50 & 304.00 \\
光伏（kW，允许调整） & 83.15 & 159.58 & 182.64 & 403.88 \\
电价（元/kWh） & 0.0511 & 0.0700 & 0.0539 & 0.0737 \\
\bottomrule\end{tabular}\end{table}
6小时误差取每次刷新后首36格，共48\,096个非重叠样本；24小时误差按发布--目标对计算，共192\,168个样本，窗口重叠并截至观测末端。允许调整条件下，光伏6小时MAE较小，24小时MAE反而较大，需分别评价短窗和长窗。两类条件同时改变光伏信息与合同权限，不能把两类费用差单独解释为官方预报带来的收益。

对其他模型机制也作了独立回放。表\ref{tab:q4-candidates}只列完整可行候选，每行仅改变所列机制，未将其联合叠加。
\begin{table}[H]\centering\small
\caption{不同独立模型候选的回放期实际总费（单位：万元）}\label{tab:q4-candidates}
\begin{tabular}{lrr}\toprule
独立机制 & 每日冻结 & 允许调整 \\ \midrule
点预测基准 & 1718.1320 & 1586.6164 \\
长窗光伏与历史基线融合 & --- & 1579.7981 \\
长窗光伏偏差校正 & --- & 1590.7360 \\
普通窗口75\%分位数残值 & --- & 1586.3447 \\
显式80\%分位数采购余量 & 1712.1735 & 1599.9834 \\
联合误差情景采购 & 1591.9083 & 1540.2971 \\
\bottomrule\end{tabular}\end{table}
表中横线表示未开展该条件的年度试验。长窗融合与残值调整在允许调整条件下略有节费，偏差校正和显式余量则增加费用；显式余量在每日冻结条件下有所改善。本轮两类条件均采用实际总费较低的联合误差情景采购。该比较基于已查看的开发年份，尚不足以证明其他年份仍有相同收益。

\subsubsection{典型日的连续运行}
图\ref{fig:q4-freeze}与图\ref{fig:q4-adjust}展示最新情景采购在12月21日的实际轨迹。各有144个动作区间和145个储能状态，日初沿用此前实际储电量。每日冻结条件下初始与最终合同重合；允许调整条件下可按发布预报更新尚未开始区间的合同。紧急购电及充放电均取实际执行量。
每日冻结条件的当日合同共98\,116.693 kWh，实际紧急购电482.034 kWh。储电量从10\,781.398变为10\,162.913 kWh。
\begin{figure}[!htbp]\centering
\PaperGraphic[width=0.94\linewidth]{figures/q4_2_2025-12-21.pdf}
\caption{情景采购在每日冻结条件下12月21日的实际运行结果}\label{fig:q4-freeze}
\end{figure}
允许调整条件有11个区间的最终合同与初始合同不同，合同总量由99\,272.136增加至99\,886.660 kWh，实际紧急购电207.906 kWh。储电量从10\,186.525变为10\,766.283 kWh。
\begin{figure}[!htbp]\centering
\PaperGraphic[width=0.94\linewidth]{figures/q4_3_2025-12-21.pdf}
\caption{情景采购在允许调整条件下12月21日的实际运行结果}\label{fig:q4-adjust}
\end{figure}
}{
  \IfFileExists{tables/q4_results.tex}{% Generated candidate Q4 evidence; explicit run IDs; no run selection approval.
\subsection{全年回放结果}
各主方案及前一日同slot价格对照从6\,000 kWh独立连续运行。四条轨迹均覆盖334日、48\,096区间。完整轨迹实际年末储能6\,000 kWh，连续SOC、合同权限、实际输入及账本复核通过。表中只在展示时舍入。
\begin{table}[htbp]\centering\small
\caption{Q4实际费用与价格预测对照（费用单位：万元）}\label{tab:q4-annual}
\begin{tabular}{llrrrr}\toprule
条件 & 价格预测 & 计划费 & 调整费 & 紧急费 & 总费 \\ \midrule
Q4-2 & 主方案 & 1296.3714 & 0.0000 & 421.7606 & 1718.1320 \\
Q4-2 & lag1 & 1300.1209 & 0.0000 & 421.3081 & 1721.4290 \\
Q4-3 & 主方案 & 1315.2016 & 45.7007 & 225.7140 & 1586.6164 \\
Q4-3 & lag1 & 1319.0238 & 46.0643 & 224.8064 & 1589.8946 \\
\bottomrule\end{tabular}\end{table}
\begin{table}[htbp]\centering\small
\caption{Q4预测误差（每格为MAE/RMSE；负载、光伏单位kW，价格单位元/kWh）}
\begin{tabular}{lllrrr}\toprule
条件 & 价格预测 & 窗口 & 负载 & 光伏 & 价格 \\ \midrule
Q4-2 & 主方案 & 6h & 163.91/227.69 & 153.40/303.86 & 0.0511/0.0700 \\
Q4-2 & 主方案 & 24h & 166.98/233.88 & 153.50/304.00 & 0.0539/0.0737 \\
Q4-2 & lag1 & 6h & 163.91/227.69 & 153.40/303.86 & 0.0843/0.1238 \\
Q4-2 & lag1 & 24h & 166.98/233.88 & 153.50/304.00 & 0.0844/0.1239 \\
Q4-3 & 主方案 & 6h & 163.91/227.69 & 83.15/159.58 & 0.0511/0.0700 \\
Q4-3 & 主方案 & 24h & 166.98/233.88 & 182.64/403.88 & 0.0539/0.0737 \\
Q4-3 & lag1 & 6h & 163.91/227.69 & 83.15/159.58 & 0.0843/0.1238 \\
Q4-3 & lag1 & 24h & 166.98/233.88 & 182.64/403.88 & 0.0844/0.1239 \\
\bottomrule\end{tabular}\end{table}
Q4-2主方案相对lag1的实际费用减少32969.31元（0.192\%），6小时价格MAE差为-0.033217元/kWh。保留这一结果，不根据全年费用重新调整预测权重。
Q4-3主方案相对lag1的实际费用减少32782.15元（0.206\%），6小时价格MAE差为-0.033217元/kWh。保留这一结果，不根据全年费用重新调整预测权重。
Q4-2主方案相对lag1紧急购电量增加1813.02 kWh。 Q4-3主方案相对lag1紧急购电量增加3026.48 kWh。 光伏MAE：Q4-3的6h/24h为83.15/182.64 kW，Q4-2为153.40/153.50 kW。短窗与长窗结论分别报告。
6小时误差取每次刷新后首36格、不重叠；24小时误差按发布与目标对计数、窗口重叠，并截到观测末端。6小时样本数为48\,096，24小时为192\,168。Q4-2和Q4-3同时改变光伏信息及合同权限，费用差不单独归因于官方预报。预测样本数、月指标和费用分项保存在绑定显式run的证据目录。
\IfFileExists{../outputs/evidence/q4_annual/q4_2_2025-12-21.pdf}{
\begin{figure}[htbp]\centering
\includegraphics[width=0.92\textwidth]{../outputs/evidence/q4_annual/q4_2_2025-12-21.pdf}
\caption{Q4-2主方案12月21日的实际量、合同、充放电和连续储能；取自同一全年轨迹。}
\end{figure}
}{}
\IfFileExists{../outputs/evidence/q4_annual/q4_3_2025-12-21.pdf}{
\begin{figure}[htbp]\centering
\includegraphics[width=0.92\textwidth]{../outputs/evidence/q4_annual/q4_3_2025-12-21.pdf}
\caption{Q4-3主方案12月21日的实际量、合同、充放电和连续储能；取自同一全年轨迹。}
\end{figure}
}{}
\noindent\textit{复现绑定：}候选运行编号为\texttt{q4-2-v3-annual-main-001, q4-2-v3-annual-lag1-001, q4-3-v3-annual-main-001, q4-3-v3-annual-lag1-001}；代码与配置快照、逐格账本及SHA256见各run目录及\texttt{outputs/evidence/q4\_annual/comparison.json}。具体运行选择与人工核验仍待参赛队确认。
}{回放结果待补充。}
}
\FloatBarrier
